%% ch02 --- The rule of the next step
%%
%% Written September 2026, on the pattern of Chapter 1. Every number the
%% reader cannot do in their head comes from code/measure/ch02.py, which
%% imports the SAME rules the listings print (code/ch02/) and chaoslab's
%% `linear` and `orbit`, which those listings use, so the page and the
%% listings cannot disagree.
%%
%% WHAT THIS CHAPTER MAY NOT SAY. It teaches iteration on STRAIGHT-LINE
%% rules, where the whole future can be read off without iterating, and it
%% introduces the slope as the factor a small error is multiplied by in one
%% step. It measures the slope of one curved rule and stops there: what an
%% orbit does when a curved rule stretches errors in some places and shrinks
%% them in others is planted as a question and not answered. Feedback and
%% the curved rule are Chapter 3's; the logistic map, and the cobweb drawn
%% on a curve, are Chapter 5's; errors that grow exponentially are
%% Chapter 6's. Nothing in this chapter is chaotic, and it never says what a
%% curved rule will do.
%%
%% Twin: chapters/pl/ch02-iteration.tex. Section for section, macro for
%% macro, number for number; tools/parity.py compares them.

\chapter{The rule of the next step}
\label{ch:iteration}
\index{iteration}\index{rule!of the next step}

\noindent
Your doctor prescribes one tablet every morning. Each tablet puts 100
milligrams of the drug into your blood, and by the next morning your body
has cleared half of whatever was there. You take one a day for a month, then
a year, then ten years. Does the amount in your blood keep climbing for as
long as you keep taking them?

Hold on to your answer; this chapter settles it. The question is about a
rule applied over and over \dash{} halve, add a tablet; halve, add a tablet
\dash{} and that one kind of question covers savings accounts, cooling tea,
growing populations, loans, and a hotel shower whose temperature you cannot
get right. Chapter~\ref{ch:models} called the line
\code{temp = next\_minute(temp)} the most important line in the book. This
chapter is about what that line does when you run it for ever, and by the
end of it you will be able to say what a whole family of rules will do in
the long run without running them at all.

\begin{outcomes}
\outcome{write a rule of the next step in the book's notation and work out
  the list of states it produces, by hand and in Python}
\outcome{recognise the straight-line rule behind savings, medicine, cooling
  and populations, and read its long-run behaviour from one number}
\outcome{answer a \enquote{how many steps until} question with a logarithm}
\outcome{find the state where a rule stands still, and decide whether the
  rule keeps it there when it is nudged}
\outcome{measure the slope of a rule at a point, and draw and read a cobweb
  diagram}
\end{outcomes}

\section{Today, tomorrow, the day after}
\label{sec:ch02-rule}
\index{state}\index{orbit}

\noindent
Put 1000 pounds into a savings account that pays 5 per cent interest a
year, and leave it alone. After one year the bank adds 5 per cent of 1000,
which is 50, and you have 1050. After the second year it adds 5 per cent of
1050, and you have \val{ch02.save.two}. After the third,
\val{ch02.save.three}. Every year the same thing happens to whatever is
there.

That is a rule of the next step, and it needs a little notation, because
you will meet a great many of them. Call the balance after $n$ years $x_n$,
read \enquote{$x$ sub $n$}. The small $n$ is a label saying which year, not
a number to multiply by: $x_0 = 1000$ is where you start, $x_1 = 1050$ is
the balance a year later, and so on. The bank's rule is then
\[ x_{n+1} = \num{1.05}\, x_n , \]
which says in symbols what the sentence said in words: next year's balance
is \num{1.05} times this year's. Because $n$ can be any whole number, this
one line describes every year at once.

Every stepping rule in this book has the same shape. There is a
\emph{state}, a number $x$ that says where things are now; a \emph{rule},
which the book calls $f$ and which stands for \enquote{whatever you do to
the state to get the next one}; and the next state, which is $f$ applied to
this one:
\[ x_{n+1} = f(x_n) . \]
When the step number does not matter, the book writes the rule with an
arrow, $x \mapsto \num{1.05}\,x$, read \enquote{$x$ goes to \num{1.05} times
$x$}. Start from $x_0$ and apply the rule once, and you have $x_1$; apply it
to $x_1$ and you have $x_2$. The whole list $x_0, x_1, x_2, \ldots$ is the
\emph{orbit} of $x_0$, and producing it \dash{} feeding each answer back in
as the next question \dash{} is \emph{iteration}.

\mermaidfig{ch02-feed-back}{Iteration. The rule turns the state now into
the next state, and the next state becomes the state now; the same rule is
used at every step.}{the loop of iteration}

\begin{think}
Take the rule $x \mapsto 2x + 1$: double, then add one. Starting from
$x_0 = 1$, what are $x_1$, $x_2$ and $x_3$?
\end{think}
\reveal{$3$, $7$ and $15$. Twice $1$ plus one is $3$; twice $3$ plus one
is $7$; twice $7$ plus one is $15$.}
There is no way to reach $x_3$ except through $x_1$ and $x_2$: each answer
is the next question. To know where an orbit is after a hundred steps you
take a hundred steps, which is tedious by hand and effortless for a
machine.

\begin{projectbox}
The book's library, \pkg{chaoslab}, has a function that takes the steps for
you. It is the loop from Chapter~\ref{ch:models} with the bookkeeping done
once:

\pyregion{code/src/chaoslab/iterate.py}{orbit}{Apply a rule $n$ times and keep every state}{lst:ch02-orbit}

\noindent
Its partner \code{linear(a, b)} builds the rules this chapter is about: it
hands back the rule $x \mapsto a\,x + b$, ready to give to \code{orbit}.
\end{projectbox}

\section{One straight line, three fates}
\label{sec:ch02-fates}
\index{linear rule}\index{rule!linear}

\noindent
The savings rule only multiplies. Many rules multiply and then add a fixed
amount, and you have met some of them already.

\begin{itemize}
\item \textbf{The tea} of Chapter~\ref{ch:models} loses a tenth of its gap
  to a room at 20 degrees each minute. From $x$ degrees it goes to
  $x - \num{0.1}(x - 20)$, which tidies up to $\num{0.9}\,x + 2$.
\item \textbf{The tablets} of this chapter's opening: halve, then add 100,
  which is $x \mapsto \num{0.5}\,x + 100$.
\item \textbf{A population} in which births outnumber deaths by a fifth
  each year, and nothing else changes: $x \mapsto \num{1.2}\,x$.
\item \textbf{A rumour} that everyone who knows it passes to one new person
  a day: the number who know doubles, $x \mapsto 2x$.
\end{itemize}

\noindent
All of them have one shape,
\[ x \mapsto a\,x + b : \]
multiply by a fixed number $a$, then add a fixed number $b$. Draw it as a
graph, with the state now along the bottom and the next state up the side,
and you get a straight line, which is why it is called a \emph{linear}
rule. The straight line describes the rule, not what happens over time;
what linear rules do over time is, as you are about to see, mostly curved.

Start four of them from the same place and let them run:

\pyregion{code/ch02/fates.py}{rules}{Four linear rules, iterated side by side from the same start}{lst:ch02-fates}

\transcript{ch02-fates}

\noindent
The first column doubles for ever. The second halves towards nothing. The
third creeps up on $10$ from below and never passes it. The fourth also
settles on $10$, but it gets there by jumping from one side to the other,
each jump smaller than the last; look closely and it is the third column
mirrored, the same distance from $10$ at every step, on alternating sides.
Those are the three fates of a linear rule: grow without bound, shrink to
nothing, or settle on one number. A negative multiplier adds a twist,
alternation, but not a fourth fate.

\plotfig{ch02-three-fates}{The three fates of a linear rule, from the same
start. One grows without bound, one dies away, and two settle on the same
number \dash{} one creeping up on it, one jumping from side to
side.}{three linear rules and their three fates}

\begin{think}
Back to the savings account. You want the balance after ten years, and you
would rather not take ten steps. What single calculation gives it?
\end{think}
\reveal{Multiply 1000 by \num{1.05} ten times over: $x_{10} =
\num{1.05}^{10} \times 1000$, which is \val{ch02.save.ten}.}
The small raised $10$ is a \emph{power}: $\num{1.05}^{10}$ means \num{1.05}
multiplied by itself ten times. When a rule only multiplies, $x \mapsto
a\,x$, its orbit has a formula, $x_n = a^n x_0$, and you can jump to any
step you like. Growth of that kind, where every step multiplies by the same
number, is called \emph{exponential}; its shrinking twin, with $a$ between
$0$ and $1$, is \emph{exponential decay}.

\section{Slow explosions, and how many steps until}
\label{sec:ch02-until}
\index{exponential growth}\index{doubling time}\index{logarithm}
\index{half-life}

\noindent
\enquote{Exponential} has become a word for \enquote{fast}. Here is a test
of that.

\begin{think}
Two savings accounts both start with 1000. Account A adds 50 every year.
Account B grows by 2 per cent a year, which is exponential growth. Which is
ahead after thirty years?
\end{think}
\reveal{Account A, with $1000 + 30 \times 50 = 2500$ against account B's
\val{ch02.acc.grows.thirty}. B does not overtake until year
\val{ch02.acc.cross}.}
If you chose B, you were reasoning from the word rather than from the rule.
Two per cent of 1000 is 20, less than A's 50, and B needs decades of
growing on its own growth before its 2 per cent is worth more than A's flat
50.

\begin{trapbox}
\textbf{\enquote{Exponential growth means fast growth.}} It means growth in
proportion to what is already there: every step multiplies by the same
number, and that number can be very close to one. At 1 per cent a year a
sum takes about \val{ch02.dbl.one.log} years to double, which is
exponential and slower than almost anything you would call fast. What
exponential growth guarantees is not speed but inevitability: however
small the rate, it overtakes any growth that adds a fixed amount,
eventually \dash{} and eventually can be a lifetime.
\end{trapbox}

\plotfig{ch02-two-accounts}{Adding a fixed amount against growing by a fixed
fraction. The exponential account is behind for most of a century, and then
it is not.}{linear against exponential growth}

\noindent
The useful question to ask of exponential growth is not \enquote{how fast?}
but \enquote{how many steps until it doubles?}, because the answer does not
depend on where you start. From 1000 or from \num{1000000}, a balance
growing at 7 per cent a year doubles in the same number of years.

\begin{think}
At 7 per cent a year, roughly how many years does a sum take to double:
about 5, about 10, or about 14?
\end{think}
\reveal{About 10. Counting whole years, 1000 first passes 2000 in year
\val{ch02.dbl.seven.count}; the precise answer is \val{ch02.dbl.seven.log}
years.}
Counting works, and there is also a button for it. The question is: how
many times must I multiply by \num{1.07} to reach 2? In symbols, for which
$n$ is $\num{1.07}^n = 2$? A \emph{logarithm} answers exactly that kind of
question. Your calculator has one marked \code{ln}, the natural logarithm,
and the one thing you need to know about it is what it does to a power:
\[ \ln (a^n) = n \ln a . \]
It turns \enquote{multiply by $a$, $n$ times} into \enquote{add $\ln a$, $n$
times}. Take the logarithm of both sides of $a^n = 2$ and you have $n \ln a
= \ln 2$, so
\[ n = \frac{\ln 2}{\ln a} . \]
In words: the number of steps is the logarithm of how far you want to go,
divided by the logarithm of one step. With $\ln 2 = \val{ch02.ln.two}$ and
$\ln \num{1.07} = \val{ch02.ln.seven}$, that is \val{ch02.dbl.seven.log}.

\pyregion{code/ch02/doubling.py}{count}{Two ways to ask how long a balance takes to double}{lst:ch02-doubling}

\pyregion{code/ch02/doubling.py}{table}{The doubling-time table}{lst:ch02-table}

\transcript{ch02-doubling}

\noindent
The counted column is the logarithm's answer rounded up to a whole step. At
2 per cent the logarithm says \val{ch02.dbl.two.log} and the count says
\val{ch02.dbl.two.count}: after 35 years the balance is
\val{ch02.dbl.two.short} times what went in, a hair short of double, and
the table has rounded the hair away. The last column is the banker's
shortcut, the \emph{rule of 70}: divide 70 by the percentage. It works
because for a small rate the logarithm of one step is close to the rate
itself ($\ln \num{1.07}$ is nearly \num{0.07}), and $100 \times \ln 2$ is
about 70.

Shrinking works the same way. A drug of which the body clears a fifth every
hour keeps \num{0.8} of itself each hour, and the time until half is left,
its \emph{half-life}, is $\ln \num{0.5} / \ln \num{0.8}$, about
\val{ch02.half.fifth} hours. The tablets of the opening, halved every day,
have a half-life of exactly one day.

\begin{lab}{l02_01_until}{How many steps until?}
Write \code{steps\_until}, which counts the steps of $x \mapsto a\,x$ from
$1$ until $x$ reaches a target, and \code{time\_until}, which asks the
logarithm the same question; the test checks that the count is always the
logarithm rounded up, for growth and for decay. Then write
\code{half\_life} for a drug that loses a fixed fraction every step.
Python's natural logarithm is \code{math.log}.
\end{lab}

\section{Where the rule stands still}
\label{sec:ch02-fixed}
\index{fixed point}

\noindent
Now the opening question. The state is the amount of drug in your blood
just after the morning tablet, in milligrams. Overnight half of it goes,
and the next tablet adds 100, so the rule is $x \mapsto \num{0.5}\,x + 100$,
starting from $x_0 = 100$.

\begin{think}
Work out the first few mornings: $100$, then half of that plus $100$, and
so on. Does the amount keep climbing for as long as you take the tablets?
\end{think}
\reveal{No. It goes $100$, $150$, $175$, \num{187.5}, \num{193.75}, and on,
closing in on $200$ without ever reaching it.}
Why $200$? Because $200$ is the amount the rule leaves alone: half of $200$
is $100$, and the tablet puts the $100$ straight back. A state that the rule
does not change is called a \emph{fixed point}, and you find one by asking
which $x$ satisfies
\[ x = f(x) . \]
For the tablets that is $x = \num{0.5}\,x + 100$. Take $\num{0.5}\,x$ from
both sides and $\num{0.5}\,x = 100$, so $x = 200$. For any linear rule the
same two lines give
\[ x^{*} = \frac{b}{1 - a} , \]
where the star marks the fixed point and $a$ must not be exactly $1$ (the
rule $x \mapsto x + b$ keeps adding and stands still nowhere). The tea's
rule, $x \mapsto \num{0.9}\,x + 2$, has its fixed point at
$2 / \num{0.1} = 20$ degrees: the temperature of the room, as it should.

\pyregion{code/ch02/drug.py}{tablets}{One tablet a morning, and the gap to the fixed point}{lst:ch02-drug}

\transcript{ch02-drug}

\noindent
Read the right-hand column. The shortfall is $100$, then $50$, then $25$: it
halves every morning, because every night the body halves whatever is
there, and the tablet adds the same $100$ however much that is. After
\val{ch02.drug.within} days you are within one milligram of $200$, and you
will never be exactly there. The amount in your blood does not climb for
ever. It levels off, and where it levels off is a fixed point.

\begin{worldbox}
Pharmacologists call this fixed point the \emph{steady state}, and the rule
of thumb taught to doctors and pharmacists is that a medicine taken at
regular intervals reaches it after four or five half-lives. The arithmetic
is the table above: after four half-lives the shortfall has halved four
times, leaving the level at \val{ch02.drug.four} per cent of the steady
state, and after five at \val{ch02.drug.five}. When a medicine is needed at
full strength straight away, a doctor may give a larger first dose, a
\emph{loading dose}, which starts the orbit on the fixed point rather than
below it.
\end{worldbox}

\section{Staying put, and running away}
\label{sec:ch02-stability}
\index{fixed point!stable}\index{fixed point!unstable}\index{stability}

\noindent
Here is a fixed point of a different character. You owe \num{10000} on a
loan that charges 1 per cent interest a month, and you repay 100 a month,
so each month the balance goes to $\num{1.01}\,x - 100$.

\begin{think}
What happens to the balance over the years? And what happens if you owe
\num{11000} instead, with the same interest and the same repayment?
\end{think}
\reveal{At \num{10000}, nothing, ever: the interest is exactly 100 and your
100 cancels it. At \num{11000} the interest is 110, your 100 does not cover
it, and after ten years you owe \val{ch02.loan.up}.}
Both balances obey the same rule, and \num{10000} is its fixed point:
$\num{1.01} \times \num{10000} - 100 = \num{10000}$. But the rule does not
pull you towards it. Start a thousand above it and the debt runs away
upwards; start a thousand below, at \num{9000}, and it runs away downwards,
to zero, paid off after \val{ch02.loan.payoff} months, about
\val{ch02.loan.years} years.

\begin{trapbox}
\textbf{\enquote{A fixed point is where the system ends up.}} It is where
the system stays if it is exactly there. Whether a system near it moves
towards it or away from it is a separate question, and the loan's fixed
point pushes every other balance away. A fixed point that draws nearby
states in is \emph{stable}; one that pushes them out is \emph{unstable}.
Finding a fixed point tells you where to look. Only its stability tells you
whether anything will be found there.
\end{trapbox}

\noindent
The difference between the tablets and the loan is one number. Take any
linear rule with a fixed point $x^{*}$ and follow the gap between a state
$x$ and $x^{*}$ through one step:
\[ f(x) - x^{*} = (a\,x + b) - (a\,x^{*} + b) = a\,(x - x^{*}) . \]
The first equals sign uses the fact that $x^{*}$ is fixed, so that
$a\,x^{*} + b$ is $x^{*}$ itself. The $b$ cancels, and what is left says:
\textbf{each step multiplies the gap by $a$}. For the tablets $a =
\num{0.5}$ and the gap halves; for the loan $a = \num{1.01}$ and the gap
grows by 1 per cent a month; for the tea $a = \num{0.9}$, which is why the
two cups of Chapter~\ref{ch:models} drew closer by a tenth every minute.

What matters is the size of $a$, written $\abs{a}$, which ignores its sign:
$\abs{\num{-0.5}} = \num{0.5}$. If $\abs{a}$ is below 1 the gap shrinks and
the fixed point is stable; if it is above 1 the gap grows and the fixed
point is unstable. The sign says only whether the gap keeps its side or
flips.

You have felt a negative $a$ in a hotel shower. You want 38 degrees, the
water is at 30, and the tap responds slowly, so you overdo it: each turn of
the tap moves the temperature by one and a half times the error, in the
opposite direction. The rule is $x \mapsto x - \num{1.5}(x - 38)$, which is
$x \mapsto \num{-0.5}\,x + 57$. From 30 you go to 42, then 36, then 39, then
\num{37.5}: too hot, too cold, too hot, each miss half the last. Jerk the
tap harder, by two and a half times the error, and $a = \num{-1.5}$: from
30 to 50, then 20, then 65, scalding and freezing by turns and worse every
time. Here is the whole story in one table.

\begin{center}
\begin{tabularx}{\linewidth}{@{}lXX@{}}
\toprule
the multiplier $a$ & a nudge away from $x^{*}$ & the orbit \\
\midrule
above $1$ & grows & runs away, on one side \\
between $0$ and $1$ & shrinks & creeps in, from one side \\
between $-1$ and $0$ & shrinks, flipping side & settles, overshooting every
time \\
below $-1$ & grows, flipping side & runs away, swinging wider \\
\bottomrule
\end{tabularx}
\end{center}

\begin{lab}{l02_02_fixed}{Where does it stop, and does it stay?}
For the rule $x \mapsto a\,x + b$, write \code{fixed\_point}, which solves
$x = a\,x + b$; \code{is\_stable}, which decides whether a nudge away from
the fixed point shrinks (you may decide it by experiment: start a little
off, take one step, compare the gaps); and \code{steps\_to\_settle}, which
counts the steps until an orbit is within a tolerance of its fixed point.
The tests use the tablets, the tea, the loan and the shower.
\end{lab}

\section{The slope: how much a nudge is multiplied}
\label{sec:ch02-slope}
\index{slope}

\noindent
The multiplier $a$ has a second name, and it is the one the rest of the
book uses. Draw the rule as a graph: the state now along the bottom, the
next state up the side. For the tablets that is the straight line
$y = \num{0.5}\,x + 100$, where $y$ is the height of the line above $x$.
Move along the bottom from 150 to 151, one unit, and the line rises from
175 to \num{175.5}, half a unit. The \emph{slope} of a graph at a point is
exactly this: \textbf{how far the output moves for each unit the input
moves}. A straight line has the same slope everywhere, and for
$x \mapsto a\,x + b$ that slope is $a$.

So the stability rule of the last section can be said in the words the rest
of the book uses: \emph{a small error is multiplied by the slope at every
step}. A slope smaller than 1 in size and errors die; larger, and they
grow.

A curve has a slope too, but it differs from point to point, as a hill is
steep in one place and nearly flat in another. You can still measure it at
a point, by nudging. Move the input a tiny amount $h$ to either side of
$x$, see how far the output moves, and divide:
\[ \text{slope at } x \approx \frac{f(x + h) - f(x - h)}{2h} . \]
The sign $\approx$ means \enquote{very nearly equal}, and in
\pkg{chaoslab}'s function \code{slope} the nudge $h$ is one millionth.
Nothing new is hidden in it: it is the rise divided by the run, for a run
so short that the curve has no room to bend.

\pyregion{code/ch02/slope.py}{measure}{Measuring the slope by nudging, on a straight line and on a curve}{lst:ch02-slope}

\transcript{ch02-slope}

\noindent
The line's slope is \num{0.5} wherever you measure it. The curved rule
$x \mapsto x \times x$ is steeper the further right you go.

\begin{think}
The square rule has slope $3$ at \num{1.5}. Suppose you meant to feed it
\num{1.5} but fed it \num{1.501}, an error of \num{0.001}. Roughly how wrong
is the output?
\end{think}
\reveal{By about \num{0.003}: the error is multiplied by the slope.
Exactly, $\num{1.501} \times \num{1.501} - \num{2.25} =
\val{ch02.nudge.exact}$.}
That is the fact behind the two cups and behind every gap in this chapter,
now on a curve: an error is multiplied by the slope at the place where it
happens to be. On a straight line the factor never changes, which is why
this whole chapter could be foreseen without iterating. On the square rule
it does: between $0$ and \num{0.5} the slope is below 1 and errors shrink,
beyond \num{0.5} it is above 1 and they grow. So here is a question this
chapter cannot answer. What does an orbit do if its rule is curved, so
that it keeps passing from places that shrink errors to places that
stretch them? A straight line never does that, and nothing you have learned
so far says. Chapter~\ref{ch:feedback} bends the line, and
Chapter~\ref{ch:logistic} follows an orbit that does exactly this.

\section{Seeing it: the cobweb diagram}
\label{sec:ch02-cobweb}
\index{cobweb diagram}

\noindent
There is a way to watch iteration happen on the rule's own graph, with a
pencil and no arithmetic. Draw the rule's line, and draw the diagonal
$y = x$, whose height is always equal to its position along the bottom.
Then:

\begin{enumerate}
\item Start at $x_0$ on the bottom axis and go straight up, or down, to the
  rule's line. The height you reach is $x_1$.
\item To use $x_1$ as the next input you need it back along the bottom. Go
  straight across to the diagonal: there, position equals height, so you
  are now standing above $x_1$.
\item Go up or down to the line again, and the height is $x_2$. Across to
  the diagonal, to the line, and so on.
\end{enumerate}

\noindent
The corners of the path are the orbit, drawn. The function \code{cobweb}
in \pkg{chaoslab} lists them, and for the tablets, starting from nothing in
the blood, the first few are these:

\pyregion{code/ch02/cobweb_steps.py}{corners}{The corners of the cobweb for the tablets}{lst:ch02-cobweb}

\transcript{ch02-cobweb}

\noindent
Every corner on the diagonal has the same number twice, and those numbers
are the orbit: $0$, $100$, $150$, $175$. The path can come to rest only
where the rule's line crosses the diagonal, because there the height of the
line equals its position, $x = f(x)$: the crossing is the fixed point, seen.

\begin{think}
The loan's line, $y = \num{1.01}\,x - 100$, is steeper than the diagonal.
Starting a little to one side of the crossing, does the path climb into the
crossing or away from it?
\end{think}
\reveal{Away. The rule multiplies the gap to the crossing by \num{1.01}, so
each step lands a little further out than the last, on whichever side it
started.}

\plotfig{ch02-cobweb}{Cobwebs on three lines that all cross the diagonal at
$200$. A slope below one in size draws a path into the crossing \dash{} a
staircase if the slope is positive, a web if it is negative \dash{} and a
slope above one draws it out.}{a cobweb on a line}

\noindent
When the slope is negative the path winds round the crossing, and it is
that web which gives the diagram its name. The three panels are the table
of the last section, drawn.

Put the chapter together and you have something rare in this book:
complete foresight. Given any rule of the form $x \mapsto a\,x + b$ and a
start, you can now say, without taking a single step, what it will do for
ever. Its fixed point is $b/(1 - a)$. The gap to that point is multiplied
by $a$ every step, so it shrinks if $\abs{a} < 1$ and grows if
$\abs{a} > 1$, keeping its side if $a$ is positive and flipping if $a$ is
negative, and a logarithm says how many steps any change will take. Run a
straight-line rule a million times and nothing in the orbit is a surprise.
That certainty is what the rest of the book takes away. Bend the line, so
that the slope differs from place to place, and a rule one line long can do
things no formula in this chapter foresees. Chapter~\ref{ch:feedback} makes
the first bend.

\begin{summarybox}
\begin{itemize}
\item A \textbf{rule of the next step} turns the state now into the next
  state, $x_{n+1} = f(x_n)$. Applying it again and again is
  \textbf{iteration}, and the list of states it produces is the
  \textbf{orbit}.
\item Savings, cooling, dosing and simple populations all follow a
  \textbf{linear rule}, $x \mapsto a\,x + b$, whose three fates are to grow
  without bound, shrink to nothing, or settle; a negative $a$ makes the
  orbit \textbf{alternate}.
\item Growth by a fixed fraction is \textbf{exponential}: proportional, and
  not necessarily fast. The \textbf{logarithm} answers \enquote{how many
  steps until}: changing by a factor $r$ takes $\ln r / \ln a$ steps, which
  gives the doubling time and the half-life.
\item A \textbf{fixed point} is a state the rule leaves unchanged; solve
  $x = f(x)$. For a linear rule it is $b/(1 - a)$.
\item Each step multiplies the gap to the fixed point by $a$. The fixed
  point is \textbf{stable} if $\abs{a} < 1$ and \textbf{unstable} if
  $\abs{a} > 1$, and nothing ends up at an unstable one.
\item The \textbf{slope} is how far the output moves per unit the input
  moves, and it is the factor a small error is multiplied by in one step. A
  line has one slope everywhere; a curve has a different one at each point.
\item The \textbf{cobweb diagram} draws an orbit on the rule's own graph.
  On a line it closes in on the crossing with the diagonal when the slope
  is below one in size, and leaves it when the slope is above.
\end{itemize}
\end{summarybox}

\canyou

\begin{checkpoint}
\item For the rule $x \mapsto 3x - 2$, starting from $x_0 = 2$, what are
  $x_1$, $x_2$ and $x_3$?
  \answerto{$4$, $10$ and $28$: three times $2$ less two is $4$, three
  times $4$ less two is $10$, and three times $10$ less two is $28$.}
\item Tea loses a fifth of its gap to a room at 20 degrees every minute.
  Write its rule in the form $x \mapsto a\,x + b$.
  \answerto{$x \mapsto x - \num{0.2}(x - 20)$, which is
  $\num{0.8}\,x + 4$: so $a = \num{0.8}$ and $b = 4$.}
\item A town grows by 3 per cent a year. Use the rule of 70 to estimate how
  long its population takes to double, and say what counting whole years
  gives.
  \answerto{About $70 / 3$, which is 23 years. Counting whole years, the
  population first passes double in year \val{ch02.dbl.three.count}; the
  logarithm says \val{ch02.dbl.three.log}.}
\item Find the fixed point of $x \mapsto \num{0.8}\,x + 4$. Is it stable?
  \answerto{$x = \num{0.8}\,x + 4$ gives $\num{0.2}\,x = 4$, so $x = 20$,
  the room's temperature. It is stable, because $\abs{\num{0.8}} < 1$.}
\item The rule $x \mapsto 3x - 2$ has a fixed point at $1$. Is it stable,
  and what does the first question show you?
  \answerto{Unstable: $a = 3$, so every step triples the gap. The orbit
  from $2$ starts one away from the fixed point and is then $3$, $9$ and
  $27$ away: $4$, $10$, $28$.}
\item The shower rule $x \mapsto \num{-0.5}\,x + 57$ has slope \num{-0.5}.
  You are 2 degrees too cold. How far off are you after one turn of the
  tap, and on which side?
  \answerto{1 degree too hot: the error is multiplied by \num{-0.5}, which
  halves it and flips its side.}
\item In a cobweb diagram, what shape does the path make when the rule's
  slope is between $-1$ and $0$, and where does it end up?
  \answerto{It winds round the point where the line crosses the diagonal,
  in a shrinking spiral of corners, and closes in on that crossing, which
  is the fixed point.}
\end{checkpoint}
